\documentclass[10pt]{article}
\usepackage[margin=1in]{geometry}
\usepackage{amsmath,amssymb,amsthm}
\usepackage{booktabs}
\usepackage{graphicx}
\usepackage{enumitem}
\usepackage{hyperref}
\setlist{itemsep=1pt,topsep=2pt}

\newtheorem{phenomenon}{Phenomenon}
\newtheorem{proposition}{Proposition}
\newtheorem{remark}{Remark}
\newtheorem{prediction}{Prediction}
\newcommand{\E}{\mathbb{E}}
\newcommand{\R}{\mathbb{R}}
\newcommand{\ip}[2]{\langle #1, #2 \rangle}
\newcommand{\rank}{\operatorname{rank}}

\title{\vspace{-2em}Why Bilinear GCIQL Wins on OGBench Cube and Not Elsewhere:\\
Pre-Registered Phenomena, Mechanisms, and Kill-Tests}
\author{}
\date{August 2026}

\begin{document}
\maketitle
\vspace{-3em}

\paragraph{Scope.}
This note answers the two project questions --- \textbf{(RQ1)} does the bilinear form work
well across environments, and if not, what makes it work so well in OGBench; \textbf{(RQ2)}
what property is the bilinear form leveraging --- by proposing one falsifiable
\emph{phenomenon} per question factor, stated mathematically, each with a concrete kill-test
implementable in this repo. Every empirical number below is from our own runs (the
288-run MW/MS campaign in \texttt{results/table.csv}; the completed cube metric-probe matrix
in \texttt{/scratch/.../factored-fb-metric-probe/}; dataset audits of the training
\texttt{npz} files). The companion note \texttt{decoding-bilinear.tex} supplies the kernel
formalism; this note supplies the causal story and its tests.

\section{The evidence to be explained}\label{sec:evidence}

Notation: critics $\in\{$\textsc{biql} (bilinear GCIQL), \textsc{gciql} (monolithic),
\textsc{crl}, \textsc{fb}, \textsc{qrl}, \textsc{vip}$\}$; actors $\in\{$\textsc{bcflow}
(flow-matching BC, $q_{\mathrm{coeff}}{=}0$: the critic \emph{never} touches the policy),
\textsc{flowbc} (identical plus the critic's $Q$-ascent term), \textsc{ddpgbc}$\}$. All
critics share HPs across environments: $\gamma=0.99$, expectile $\tau=0.9$, lr $3{\cdot}10^{-4}$,
batch 1024, hidden $[512,512,512]$ + LayerNorm, goal sampling
$(p_{\mathrm{cur}},p_{\mathrm{traj}},p_{\mathrm{rand}})=(0.2,0.5,0.3)$ geometric.

\begin{itemize}
\item[\textbf{E1}] \textbf{Cube (OGBench \texttt{cube-single-play-v0}, 500k, 3 seeds).}
  \textsc{biql}-$d50$+\textsc{flowbc} $=92.7\pm1.4$ (fresh reproduction of the historical
  $90.7$); \textsc{biql}-$d512=89.2$; \textsc{gciql}$=73.2\pm1.3$; \textsc{crl}$=28$--$34$;
  pure-BC floor (\textsc{bcflow} arms) $=8$--$10$. Thus on cube the \emph{entire} signal is
  critic guidance: $+84$ points for \textsc{biql}, $+64$ for \textsc{gciql}, $+20$ for
  \textsc{crl} over the same BC policy.
\item[\textbf{E2}] \textbf{Metaworld.} The BC floor is already the ceiling:
  \textsc{bcflow} alone scores $0.89$--$1.00$ on play-noise (\texttt{mw-push},
  \texttt{mw-pick-place}, historical \texttt{mw-bin-picking}, \texttt{mw-button-press}) and
  $0.58$--$0.93$ on expert. Adding the critic ($\to$\textsc{flowbc}) changes success by
  $-0.27$ to $+0.00$; \textsc{biql}${\approx}$\textsc{gciql} everywhere.
\item[\textbf{E3}] \textbf{ManiSkill (play-noise).} Critic guidance has \emph{negative} sign
  for every goal-conditioned critic and \emph{positive} sign for \textsc{fb}:
  on \texttt{ms-pick-cube}, \textsc{bcflow}$\to$\textsc{flowbc} moves \textsc{biql}
  $0.38\to0.08$, \textsc{gciql} $0.42\to0.14$, \textsc{crl} $0.40\to0.10$, but \textsc{fb}
  $0.57\to0.69$.
\item[\textbf{E4}] \textbf{The bilinear critic does \emph{not} have the best temporal
  geometry.} Probe P1 (Spearman of $\hat d$ with time-gap, partialed on raw-obs distance):
  \textsc{gciql} $0.63/0.54$ $>$ \textsc{biql} $0.54/0.41$. RSA across methods:
  \textsc{gciql}$\leftrightarrow$\textsc{qrl} $=0.88$,
  \textsc{gciql}$\leftrightarrow$\textsc{crl} $=0.54$ (a three-method consensus geometry),
  while \textsc{biql} agrees with nobody ($\le 0.21$) and is closest to \emph{raw
  observation distance} ($0.45$--$0.52$). Yet \textsc{biql} is the best policy by 20 points.
\item[\textbf{E5}] \textbf{The bilinear solution is a bounded, norm-concentrated, low-rank
  object.} $\mathrm{corr}(\ip{\varphi}{\psi},\cos\angle(\varphi,\psi))=0.996$;
  norm CVs $\le 0.09$ (both towers, both $d$); effective rank $28/50$ and $48/512$;
  $d{=}5$ historically collapses to $9.6\%$. (Caveat: the untrained net already has
  cos-corr $0.994$ --- norm concentration is architectural, from LayerNorm; what is
  \emph{learned} is the direction field.)
\item[\textbf{E6}] \textbf{Geometry saturates early; success keeps climbing.} From 100k to
  500k steps, P1-partial moves $0.37\to0.41$ while success moves $65\to93$.
\item[\textbf{E7}] \textbf{The datasets are not the same kind of object.}
  Cube-play: $10^6$ transitions, 1000 trajectories of length 1000, task-agnostic scripted
  play that relocates the block indefinitely (no terminal dwell). MW: $10^5$ transitions,
  length 100, (noisy-)expert reaching one target then \emph{dwelling at the goal for
  48--71\% of all recorded steps}. MS: $2.5{\times}10^4$ transitions, length 25, median
  first-success at step 5 (expert) / 11 (play-noise), dwell 27--75\%. ``Play-noise'' for
  MW/MS is $\varepsilon$-mixed noisy \emph{expert} ($\varepsilon_{\max}=0.25$,
  $\sigma_{\max}=0.15$, 30\% clean) --- single-corridor data, not play.
\item[\textbf{E8}] \textbf{The goal channel is structurally different.} MW/MS are
  goal-\emph{observable}: dims $36{:}39$ (MW) / $26{:}28$ (MS) of the observation are the
  episode's target coordinates, constant within an episode (std $<10^{-6}$), redrawn across
  episodes. OGBench cube observations contain \emph{no} goal information. Consequently on
  MW/MS: (i) the $p_{\mathrm{rand}}=0.3$ cross-trajectory relabeled pairs pair a state whose
  embedded target is $c$ with a goal whose embedded target is $c'\ne c$ --- physically
  unreachable pairs; (ii) at eval, \texttt{env/unified.py:167} conditions every episode on
  a \emph{single fixed dataset frame} (argmax-reward next-obs) while the environment redraws
  its target every reset (\texttt{\_freeze\_rand\_vec=False}); (iii) \textsc{fb} is immune
  because its $z=\E[r\,B(s')]$ is inferred from the recorded \emph{task} rewards, not from a
  goal frame.
\end{itemize}

\section{Setup}

Let observations decompose as $x=(s,c)$ where $s$ is the controllable state and $c$ an
episode-immutable context (the embedded target; on cube, $c=\varnothing$). Let
$d^*(x,x')$ be the minimal dataset-supported hitting time and
$K^*=\gamma^{d^*}\in[0,1]$ the discounted reachability kernel
(\texttt{decoding-bilinear.tex}, Eq.~2). A goal-conditioned critic trained on relabeled
pairs estimates $\hat K\approx K^*$ on the sampled pair distribution $\mu$. The one-step
actor solves
\begin{equation}
\max_\theta\; J(\theta)=\E_{s,g}\Big[\,q\cdot \hat Q\big(s,\pi_\theta(s,g),g\big)\;-\;
\beta\,\|\pi_\theta(s,g)-\pi_{\mathrm{BC}}(s,g)\|^2\Big],
\label{eq:actor}
\end{equation}
with $q\in\{0,1\}$ the \textsc{bcflow}/\textsc{flowbc} switch. Everything below is about
when and why the $\hat Q$-term in \eqref{eq:actor} helps.

\section{RQ1: why the cube result does not replicate}

\subsection{Phenomenon I: guidance headroom censoring}

\begin{phenomenon}[Headroom censoring]\label{ph:headroom}
Write $S_q$ for success of the actor trained with critic weight $q$ in \eqref{eq:actor}.
Then
\begin{equation}
S_1 \;=\; S_0 \;+\; H\cdot \rho, \qquad
H := S_{\max}-S_0 \;\;(\text{headroom}), \quad \rho\in[-S_0/H,\,1]\;\;(\text{guidance quality}).
\end{equation}
The benchmark \emph{measures} $H\rho$, not $\rho$. On cube $H\approx0.9$ and the observed
$H\rho$ separates critics by 60+ points (E1). On MW $H\le0.1$ (E2): \emph{any} critic ---
perfect or useless --- produces $|S_1-S_0|\le 0.1$. The MW experiments are censored, not
null: they carry almost no information about $\rho$, hence none about the bilinear form.
\end{phenomenon}

The censoring is caused by the data, not the simulator: noisy-expert corridors (E7) make
$\pi_{\mathrm{BC}}$ already near-optimal for the eval task family. Cube-play is the opposite
regime --- BC on undirected play is at $9\%$, so \emph{all} performance must come through
$\hat Q$.

\begin{prediction}\label{pred:headroom}
Reopen headroom on MW and the ranking must reappear if and only if $\rho$ differs across
critics. Implementation (cheap, data-side only): (i) subsample MW play-noise to
$\{25\mathrm{k}, 50\mathrm{k}\}$ transitions and/or drop the 30\% clean-expert fraction and
raise $\varepsilon_{\max}$ at collection; (ii) rerun the $\{$\textsc{biql},
\textsc{gciql}$\}\times\{$\textsc{bcflow}, \textsc{flowbc}$\}$ square, 3 seeds. Falsifier of
Ph.~\ref{ph:headroom}: BC floor drops (headroom opens) but all $S_1-S_0$ remain $\approx0$
--- then the failure is not censoring but guidance invalidity (Ph.~\ref{ph:context}).
\end{prediction}

\subsection{Phenomenon II: contextual goal misspecification}

\begin{phenomenon}[Ill-posed goal channel on goal-observable benchmarks]\label{ph:context}
With $x=(s,c)$ and context immutable within an episode, the true kernel factorizes as
\begin{equation}
d^*\big((s,c),(g,c')\big)=\begin{cases} d^*_c(s,g) & c'=c\\ \infty & c'\ne c\end{cases}
\qquad\Longrightarrow\qquad
K^*\big((s,c),(g,c')\big)=\mathbf{1}[c{=}c']\,\gamma^{d^*_c(s,g)}.
\label{eq:context}
\end{equation}
Under relabeling with $p_{\mathrm{rand}}=0.3$, 30\% of value-training pairs sit in the
$c\ne c'$ block, whose regression target is the constant $-1/(1-\gamma)$; the critic must
spend capacity learning the diagonal-block indicator instead of geometry. At eval the agent
is conditioned on one fixed frame $g_0=(\bar g,c_0)$ while the environment draws
$c\sim P(c)$ fresh each episode, so almost surely $c\ne c_0$ and the conditioned pair lies
in the unreachable block: the ascent direction
$\nabla_a \hat Q\big((s,c),a,(\bar g,c_0)\big)$ is a gradient of a function whose true value
is constant, i.e.\ pure critic-error. Its expected contribution to \eqref{eq:actor} is
adversarial noise on top of BC, so $\rho<0$: guidance \emph{hurts}, exactly the sign
observed on MS (E3), while pure BC survives by reading the true target from its own
observation. \textsc{fb}'s task-inferred $z$ never touches this channel --- explaining its
immunity (E3) --- so the MW/MS tables compare a well-posed conditioning (\textsc{fb})
against an ill-posed one (all GC critics); they are not a test of the bilinear form.
\end{phenomenon}

\begin{prediction}\label{pred:context}
Three independent repairs, each implementable in a day, each with a pre-registered sign:
(a) \textbf{mask the context dims} out of the goal input ($\psi(g)$ and actor conditioning)
--- removes the $c$-block structure from $\hat K$;
(b) \textbf{per-episode goal synthesis at eval}: build $g$ from the \emph{current} episode's
observation by writing the target coordinates into the object slots (both MW and MS expose
them), instead of the fixed dataset frame;
(c) \textbf{$p_{\mathrm{rand}}=0$} for MW/MS --- removes the unreachable block from training.
Ph.~\ref{ph:context} predicts (b) alone flips the guidance sign on \texttt{ms-pick-cube}
($S_1-S_0$ from $-0.30$ toward $\ge 0$), and (a)+(c) increase it further. Falsifier: repairs
applied, headroom present (MS has $H\approx0.6$), and GC guidance still negative --- then
the defect is in the value target itself (Ph.~\ref{ph:kernel}).
\end{prediction}

\subsection{Phenomenon III: kernel degeneracy under short-horizon dwell data}

\begin{phenomenon}[Range compression and goal-marginal collapse]\label{ph:kernel}
The information the critic can add over BC lives in the \emph{variation} of $K^*$ across
the sampled pairs. On trajectories of length $H_{\mathrm{traj}}$, relabeled hitting times
satisfy $d^*\le H_{\mathrm{traj}}$, so
\begin{equation}
K^*\in[\gamma^{H_{\mathrm{traj}}},1]:\qquad
\underbrace{1-\gamma^{1000}\approx 1.0}_{\text{cube}}\qquad
\underbrace{1-\gamma^{100}=0.63}_{\text{MW}}\qquad
\underbrace{1-\gamma^{25}=0.22}_{\text{MS}} .
\end{equation}
On MS, additionally, $(1-\gamma)d^*\le0.25$ puts $K^*$ in its linear regime,
$K^*\approx 1-(1-\gamma)d^*$, so the kernel matrix is affine in $d^*$
($K\approx \mathbf{1}\mathbf{1}^\top-(1-\gamma)D$): the exponential structure whose
low-rank factorization is \textsc{biql}'s inductive bias (\texttt{decoding-bilinear.tex},
\S2) degenerates, and monolithic and bilinear parametrizations face the same easy
regression squeezed into a $0.22$-wide range against TD noise. Dwell (E7) compounds this:
48--75\% of relabeled goals collapse into the tiny at-goal set, so the goal marginal is
nearly a point mass and most pairs have $d^*\approx$ const. In this regime no critic can
differentiate actions ($\rho\approx0$ for all), and the bilinear form has nothing to
leverage --- \emph{independently} of Ph.~\ref{ph:context}.
\end{phenomenon}

\begin{prediction}\label{pred:kernel}
(a) \textbf{HP-side} (cheapest possible test, pure config): set the discount per-environment
by the scaling rule $\gamma_{\mathrm{env}}=\exp(-1/\bar d)$ with $\bar d$ the median
relabeled hitting time ($\bar d\!\approx\!10$ on MS $\Rightarrow\gamma\approx0.9$;
$\bar d\!\approx\!30$ on MW $\Rightarrow\gamma\approx0.97$), restoring the cube-like ratio
$K$-range $\approx1$. Prediction: GC critic guidance on repaired-goal MS (after
Pred.~\ref{pred:context}) becomes positive and \textsc{biql} separates from \textsc{gciql}.
(b) \textbf{Data-side} (the decisive test): collect OGBench-style \emph{play} on MS/MW ---
long horizon ($\ge500$ steps), scripted policy that chains random object relocations,
no dwell, no per-episode target reset if the backend allows freezing
(\texttt{\_freeze\_rand\_vec=True} matched between collection and eval). This equalizes E7
across benchmarks. Prediction: with E7 and E8 equalized, the cube ordering
(\textsc{biql} $>$ \textsc{gciql} $\gg$ BC) reappears on MW/MS. This is the single
experiment that most directly answers RQ1.
\end{prediction}

\section{RQ2: what the bilinear form is actually leveraging}

\subsection{Phenomenon IV: the advantage is not a better metric (established negative)}

\begin{phenomenon}[Metric fidelity does not explain guidance value]\label{ph:notmetric}
If \textsc{biql}'s edge came from a more faithful temporal geometry, it would beat
\textsc{gciql} on the geometry probes and agree with the architecturally-constrained
reference (\textsc{qrl}). The data show the opposite (E4): \textsc{gciql} tracks temporal
distance better, joins the \textsc{qrl}/\textsc{crl} consensus (RSA $0.88/0.54$), and
\emph{loses by 20 points}; \textsc{biql} is a geometric outlier nearest raw-obs distance
and wins. Moreover geometry saturates at 100k while success keeps rising (E6). Hence
value-geometry fidelity on the data distribution is neither necessary nor sufficient for
guidance value. Whatever the bilinear form buys, it buys it \emph{at the actor interface},
in properties of $\hat Q$ invisible to $(s,g)$-geometry probes.
\end{phenomenon}

This also disposes of the naive reading of H-sim: yes, the critic is numerically a scaled
cosine (E5), but so is the untrained net; cosine-ness per se is architectural, not the
learned asset.

\subsection{Phenomenon V: bounded low-rank guidance (the proposed mechanism)}

\begin{phenomenon}[Bounded, goal-uniform, low-rank action guidance]\label{ph:mechanism}
Decompose the actor's gradient signal. For the bilinear critic
$\hat Q(s,a,g)=\ip{\varphi(s,a)}{\psi(g)}/\sqrt d$ with norm concentration
$\|\varphi\|\approx\kappa_\varphi$, $\|\psi\|\approx\kappa_\psi$ (E5: CV $\le0.09$):
\begin{align}
\text{(boundedness)}\qquad & |\hat Q| \;\le\; \kappa_\varphi\kappa_\psi/\sqrt d \;=:\;\kappa
&&\text{uniformly in } (s,a,g),\label{eq:bound}\\
\text{(goal-uniform scale)}\qquad & \nabla_a \hat Q \;=\; J_a\varphi(s,a)^\top\,\psi(g)/\sqrt d,
\quad \|\nabla_a \hat Q\|\le \|J_a\varphi\|_{\mathrm{op}}\,\kappa_\psi/\sqrt d,\label{eq:grad}\\
\text{(rank restriction)}\qquad & \nabla_a\hat Q \in
\operatorname{row}(J_a\varphi)\;\text{with } \hat Q \text{ of rank} \le d\ (\text{eff.\ rank }\approx 28).\notag
\end{align}
Consequences, each targeting a known failure mode of critic-guided actors:
\begin{enumerate}[label=(\alph*)]
\item \textbf{Exploitation bound.} The actor ascends $\hat Q$, and its true regret relative
to BC obeys the standard decomposition
$J^{\pi}\ge J^{\pi_{\mathrm{BC}}}+\Delta_{\hat Q}-2\sup_{a}|\hat Q-Q^*|$ over reached
actions. For a monolithic MLP the sup is uncontrolled off-data --- the actor manufactures
phantom value --- while \eqref{eq:bound} caps what ascent can ever collect at $2\kappa$.
The observed \textsc{gciql} triangle-violation \emph{magnitude} of $9.5$ steps versus
\textsc{biql}'s $4.8$ (at comparable rates) is exactly the signature: monolithic errors are
rarer but huge; bilinear errors are small because the hypothesis class cannot express large
outliers around a rank-28 solution.
\item \textbf{Goal-uniform guidance scale.} \textsc{flowbc} normalizes the $Q$-term by the
\emph{global} mean $|\hat Q|$; per-goal scale heterogeneity survives that normalization.
By \eqref{eq:grad} with $\kappa_\psi$ concentrated (CV $0.07$), the bilinear critic feeds
every goal's gradient at the same magnitude, so no goal region dominates the actor's
batch gradient. The monolithic critic offers no such guarantee.
\item \textbf{Matrix completion to novel eval pairs.} Eval pairs $(s,g^{\mathrm{task}})$
are combinations never co-observed in play data. A rank-$d$ factorization is determined by
marginal coverage --- each state seen with \emph{some} goals, each goal with \emph{some}
states --- and fills in unseen blocks of $K^*$ (matrix-completion generalization), which is
precisely where an unstructured MLP must extrapolate freely. This is why \textsc{biql} can
look \emph{worse} on-distribution (E4 probes are val-set pairs) yet act better on eval
pairs.
\end{enumerate}
Finally, the mechanism is a \textbf{conjunction}, which our own 500k matrix already
factorizes (E1): the guidance target must be in-sample \emph{optimal} (expectile), and the
parametrization must be bounded-low-rank:
\begin{center}
\begin{tabular}{lcc}
\toprule
 & monolithic & bilinear \\
\midrule
behavior-policy target (NCE) & --- & \textsc{crl}: $28$--$34\%$ \\
in-sample-optimal target ($\tau{=}0.9$) & \textsc{gciql}: $73\%$ & \textsc{biql}: $\mathbf{93\%}$ \\
\bottomrule
\end{tabular}
\end{center}
(BC floor $9\%$.) \textsc{crl} shares the bilinear geometry and norm concentration (E5
holds for it too) but ranks actions by \emph{behavior} occupancy, so its guidance is capped;
\textsc{gciql} shares the optimal target but not the bounded-low-rank guidance. Only the
conjunction clears $90\%$.
\end{phenomenon}

\begin{prediction}[Mechanism kill-tests, all on cube, all single-file changes]\label{pred:mech}
Each sub-claim gets its own control; together they separate (a), (b), (c):
\begin{enumerate}
\item \textbf{Bounded monolithic} (tests (a) alone): train \textsc{gciql} with output
squashing $\hat Q=\kappa\tanh(\mathrm{MLP}/\kappa)$, $\kappa$ set to the observed bilinear
$\kappa\approx|V|_{\max}$, or spectral-norm the last layer. Prediction: recovers a large
fraction of the $73\to93$ gap. If it recovers \emph{none}, (a) is falsified and (c) carries
the phenomenon.
\item \textbf{De-concentrated bilinear} (tests (a)+(b) from the other side): remove
LayerNorm from the $\psi$ tower / add a norm-spread penalty so CV $>0.5$. Prediction:
success drops toward \textsc{gciql} despite identical rank.
\item \textbf{Eval-pair value probe} (tests (c), pure analysis on existing checkpoints):
run probe P2 (goal-approach monotonicity) along \emph{eval-task} rollouts --- the deferred
probe --- for \textsc{biql} vs \textsc{gciql}. Prediction: \textsc{gciql}'s ranking quality
degrades from val-pairs to eval-pairs strictly more than \textsc{biql}'s (an interaction
effect, not a main effect).
\item \textbf{Rank threshold} (tests (c)): $d\in\{5,10,20,30,50\}$, 3 seeds. Prediction:
success transitions sharply at $d\approx$ effective rank of the solution ($\approx28$),
i.e.\ between $d{=}20$ and $d{=}30$, not gradually.
\item \textbf{Fourth cell of the $2\times2$}: give \textsc{crl}'s bilinear critic an
in-sample-optimal target (expectile-weighted NCE or TD-InfoNCE). Prediction: it joins
\textsc{biql} near $90\%$; conversely \textsc{biql} with $\tau=0.5$ (SARSA-like) should
fall toward \textsc{crl}.
\end{enumerate}
\end{prediction}

\begin{remark}[Post-registration result: the $d$-sweep (P\ref{pred:mech}.4) has run]
\label{rem:dsweep}
3 seeds $\times$ 500k on cube ({\footnotesize\texttt{/scratch/.../factored-fb-metric-probe/biql-d\{5,10,20,30,50\}}}):
$d{=}5$: $0.06\pm0.05$; $d{=}10$: $0.62\pm0.07$; $d{=}20$: $\mathbf{0.92\pm0.02}$;
$d{=}30$: $0.90\pm0.02$; $d{=}50$: $0.93\pm0.02$. The transition is sharp --- $+56$ points
from $d{=}5$ to $d{=}10$, $+30$ from $d{=}10$ to $d{=}20$, then flat --- but it sits
\emph{below} the $d{=}50$ solution's effective rank of $28$: $d{=}20$ is already
statistically indistinguishable from $d{=}50$. Verdict on Pred.~\ref{pred:mech}.4: the
strong reading (``success transitions at $d\approx$ eff-rank of the unconstrained
solution'') is \textbf{falsified}; the weak reading survives and sharpens: the guidance
signal is a low-rank object of intrinsic dimension $\approx 10$--$20$, and the eff-rank
$28$ measured at $d{=}50$ \emph{overshoots} the minimal sufficient rank rather than
marking a threshold. Two follow-ups: (i) run the probe battery on the $d\in\{10,20,30\}$
checkpoints --- if learned eff-rank saturates near $15$--$20$ regardless of $d$, the
solution rank (not the container size) is the invariant; (ii) $d{=}15$ fills the one
remaining gap in the curve.

Follow-up (i) has now run (3 seeds each): learned eff-rank (phi/psi means) is
$4.1/8.0/14.7/19.6/28.0$ at $d=5/10/20/30/50$ (and $48$ at $512$) --- it does \emph{not}
saturate at an invariant solution rank; it tracks the container sublinearly (compression
ratio $0.82\to0.56\to0.09$ as $d$ grows). So eff-rank measures container-conditioned
spread, not the signal's intrinsic dimension; the only $d$-invariant quantity in the sweep
is the \emph{success saturation point} $d^\dagger\!\in\![10,20]$. Rank is a capacity
budget, not a learned constant.
\end{remark}

\begin{remark}[Post-registration result: P\ref{pred:mech}.3 has run --- the interaction is
confirmed, in a sharper form]\label{rem:p2}
Full $2\times2$ (scoring critic $\times$ rollout-generating policy), 3 seeds, 50
temperature-0 eval episodes per run, per-episode Spearman $\rho(\hat V(s_t,g),t)$
(median over episodes, mean$\pm$std over seeds; bump rate $=\Pr[\hat V_{t+1}<\hat V_t]$):
\begin{center}
\begin{tabular}{lcccc}
\toprule
scored on $\downarrow$ / rollouts from & \textsc{biql} policy ($0.95$ succ.) & \textsc{gciql} policy ($0.73$ succ.)\\
\midrule
\textsc{biql} critic  & $0.983\pm0.003$ \;(bump $0.17$) & $0.871\pm0.060$ \;(bump $0.27$)\\
\textsc{gciql} critic & $\mathbf{0.355\pm0.239}$ \;(bump $0.38$) & $0.883\pm0.015$ \;(bump $0.27$)\\
\bottomrule
\end{tabular}
\end{center}
Pre-registered prediction (gciql degrades from val-pairs to eval-pairs more than biql):
\textbf{confirmed} --- on val-set pairs \textsc{gciql} ranks temporal distance \emph{better}
(P1: $0.63$ vs $0.54$), on eval-task rollouts of the successful policy its value signal is
noise ($0.355\pm0.239$, with per-task medians as low as $-0.80$) while \textsc{biql} is a
near-perfect monotone ramp ($0.983$). The sharper form is the \emph{transfer asymmetry}:
each critic is coherent along its own policy's trajectories (diagonal $0.98/0.88$), but
only the bilinear critic transfers to the \emph{other} policy's state distribution
($0.87$ vs $0.36$). The monolithic critic's ranking is intact exactly up to the states its
own (worse) policy reaches and breaks on the states the better policy actually visits ---
which is what Ph.~\ref{ph:mechanism}(c) (completion to unseen $(s,g)$ blocks) predicts, and
what explains how ascent on $\hat Q_{\textsc{biql}}$ can push the policy \emph{beyond} the
critic-fidelity frontier of $\hat Q_{\textsc{gciql}}$ ($73\to93$).
\end{remark}

\section{Interpretation discipline}

The five phenomena are deliberately stated so that the six experiments in
Predictions~\ref{pred:headroom}--\ref{pred:mech} form a decision tree, not a shopping list:
\begin{center}
\small
\begin{tabular}{lll}
\toprule
Experiment & Outcome A $\Rightarrow$ & Outcome B $\Rightarrow$\\
\midrule
P\ref{pred:headroom} shrink MW data & ranking appears: Ph.~\ref{ph:headroom} confirmed &
flat: escalate to Ph.~\ref{ph:context}\\
P\ref{pred:context}b eval-goal synthesis (MS) & guidance sign flips: Ph.~\ref{ph:context} &
still negative: Ph.~\ref{ph:kernel} (target broken)\\
P\ref{pred:kernel}b OGBench-style play on MS & cube ordering reappears: RQ1 closed &
does not: simulator-intrinsic residual, profile dynamics\\
P\ref{pred:mech}.1 bounded \textsc{gciql} & gap closes: boundedness is the asset &
gap persists: completion (c) is the asset\\
P\ref{pred:mech}.4 $d$-sweep & sharp at $d\!\approx\!28$: rank story exact &
gradual: rank is proxy, revisit\\
\bottomrule
\end{tabular}
\end{center}

Priority order under the 12-hour budget: \textbf{P\ref{pred:mech}.3} (hours; existing
checkpoints, analysis only), \textbf{P\ref{pred:mech}.1--2} (one config/one-file edit each,
500k cube runs), \textbf{P\ref{pred:context}b} (eval-side only, no retraining),
then P\ref{pred:headroom}, P\ref{pred:kernel}. The first three alone are sufficient to
either confirm or kill Ph.~\ref{ph:mechanism}, which is the paper's central claim.

\paragraph{Honest caveats.} (i) Ph.~\ref{ph:context} explains the \emph{sign} on MS and the
FB asymmetry but not why MW-play-noise \textsc{flowbc} sometimes stays high; censoring
(Ph.~\ref{ph:headroom}) covers that, and the two are confounded until
Pred.~\ref{pred:headroom} runs. (ii) The exploitation bound in Ph.~\ref{ph:mechanism}(a) is
an upper-bound argument, not an identity; the kill-test is P\ref{pred:mech}.1, not the
algebra. (iii) VIP rows are uninformative in the current matrix: the embedding collapsed to
effective rank $1.0$ (its P1 $\approx0.04$); fix its objective/normalization before citing
it as a reference point.

\end{document}
